\clearpage
\SourceLeaf{470}


\Needspace{9\baselineskip}\begin{center}ADDITIONS.\end{center}
\BookRule{0.15\linewidth}
\textit{On Article} 28. The indeterminate equation \(ax=by\pm1\) was first solved not by the distinguished Euler (as stated there), but already by Bachet de Meziriac, a geometer of the \(17\)th century and the celebrated editor and commentator of Diophantus, for whom the distinguished La Grange claimed this honor (\textit{Add. à l'Algèbre d'Euler p.} 525, where the nature of the method is also indicated). Bachet gave his discovery in the second edition of his book \textit{Problèmes plaisans} \textit{et délectables qui se font par les nombres}, 1624; it does not yet appear in the first edition (Lyon 1612), which alone I have been able to see, although it is already announced there.

\textit{On Articles} 151, 296, 297. The distinguished Le Gendre set out his proof again in his excellent work \textit{Essai d'une théorie des nombres}, p. 214 ff., but without any essential change; this method therefore remains subject to all the objections raised in Article 297. The theorem (on which one assumption rests) that primes occur in every arithmetic progression \(l,\allowbreak l+k,\allowbreak l+2k\), etc., if \(k\) and \(l\) have no common divisor, has indeed been considered at greater length in this work, p. 12 ff.; but the requirements of geometric rigor still seem not to have been met. Yet even when this theorem has been fully proved, the other assumption will remain (that there are primes of the form \(4n+3\) for which a given prime of the form \(4n+1\), taken positively, is a quadratic nonresidue), and I do not know whether this can be proved \textit{rigorously} unless the fundamental theorem itself is already \textit{assumed}. Further, it should be observed that the distinguished Le Gendre did not make this latter assumption tacitly, but himself also acknowledged it, p. 221.
\SourceLeaf{471}

\textit{On Articles} 288---293. On the same subject, presented here as a special application of the theory of ternary forms, and apparently so fully treated in rigor and generality that nothing further could be desired, the distinguished Le Gendre undertook a much more extensive investigation in Part III of his work, pp. 321---400.\HistoricalNote{*)}{Even without our warning, readers will take care not to confuse our ternary forms with what the distinguished Le Gendre called \textit{forme trinaire d'un nombre}. By this expression he meant the decomposition of a number into three squares.} He used principles and methods entirely different from ours, but on this path became involved in several difficulties that prevented him from supporting the principal theorems by a rigorous proof. He himself candidly indicated these difficulties; but, unless we are mistaken, these may perhaps be removed more easily than the difficulty that here too the theorem just mentioned (In every arithmetic progression, etc.) is assumed, p. 371, end of the note.

\textit{On Article} 306 VIII. In the third thousand of negative determinants, \(37\) irregular ones have been found, of which \(18\) have index of irregularity \(2\), and the remaining \(19\) have index \(3\).

\clearpage\textit{On the same article} X. The question posed here has recently been completely solved. We will present this investigation, which casts remarkable light on several parts of both higher arithmetic and analysis, in the continuation of this work as soon as possible. It has also shown that the coefficient \(m\) in Article 304 is \(=\gamma\pi=2{,}3458847616\), where \(\gamma\) denotes the same quantity as in Article 302, and \(\pi\), as there, the semicircumference of a circle whose radius is \(1\).
\BookRule{0.28\linewidth}


